\documentclass[../../main.tex]{subfiles}

\begin{document}

\chapter{Integritas Data: Fungsi Hash dan MAC}
\label{chap:hash-mac}

\begin{learningoutcome}
Setelah mempelajari bab ini, mahasiswa diharapkan mampu:
\begin{itemize}
    \item Menjelaskan konsep dasar fungsi hash dan perbedaannya dengan enkripsi (Sub-CPMK 1.4).
    \item Menganalisis properti utama fungsi hash kriptografis, termasuk resistansi terhadap pre-image, second pre-image, dan kolisi (Sub-CPMK 1.4).
    \item Menguraikan mekanisme kerja algoritma hash populer seperti MD5 dan keluarga SHA (Secure Hash Algorithm) (Sub-CPMK 1.4).
    \item Menjelaskan konsep Message Authentication Code (MAC) dan HMAC untuk menjamin integritas dan otentikasi pesan (Sub-CPMK 1.4).
    \item Mengevaluasi risiko serangan kolisi pada fungsi hash dan dampaknya terhadap keamanan sistem (Sub-CPMK 1.4).
\end{itemize}
\end{learningoutcome}

\subfile{section-01}
\subfile{section-02}
\subfile{section-03}
\section{Contoh Terhitung}
Setiap angka pada contoh berikut diverifikasi secara numerik dengan Python.

\subsection{Biaya Serangan Birthday pada MD5 dan SHA-256}
Untuk peluang kolisi sekitar $50\%$, jumlah pesan acak yang harus dicoba adalah $k \approx 1{,}177 \cdot 2^{n/2}$.

\textbf{(a) MD5}, $n = 128$:
\[ k \approx 1{,}177 \cdot 2^{64} = 1{,}177 \cdot 18.446.744.073.709.551.616 \approx 2{,}17\times10^{19}. \]

\textbf{(b) SHA-256}, $n = 256$:
\[ k \approx 1{,}177 \cdot 2^{128} = 1{,}177 \cdot (3{,}402823669\times10^{38}) \approx 4{,}01\times10^{38}. \]

Rasio keduanya adalah $2^{128}/2^{64} = 2^{64} = 18.446.744.073.709.551.616$. Jadi SHA-256 membutuhkan sekitar $1{,}8\times10^{19}$ kali lebih banyak usaha daripada MD5.

\subsection{Padding Merkle--Damgård}
Aturan padding Merkle--Damgård untuk blok 512 bit: tambahkan $1$ bit ``1'', lalu $k$ bit ``0'', lalu panjang pesan $L$ dalam 64 bit, dengan syarat
\[ L + 1 + k + 64 \equiv 0 \pmod{512}, \qquad 0 \le k < 512 . \]

\textbf{Contoh: pesan \code{"abc"}.} Panjang $L = 3 \text{ byte} = 24$ bit.
\[ 24 + 1 + k + 64 \equiv 0 \pmod{512} \;\Rightarrow\; k = 423 . \]
Total $= 24 + 1 + 423 + 64 = 512$ bit, sehingga terbentuk \textbf{1 blok}.

\textbf{Contoh: pesan $L = 1000$ bit.}
\[ 1000 + 1 + k + 64 \equiv 0 \pmod{512} \;\Rightarrow\; k = 471, \]
total $= 1000 + 1 + 471 + 64 = 1536$ bit, yaitu \textbf{3 blok}.

\subsection{HMAC dengan Fungsi Hash Mainan}
Agar struktur ipad/opad terlihat, pakai fungsi hash mainan $H(b_1,\dots,b_t) = \left(\sum_i b_i\right) \bmod 256$ dan kunci 1 byte $K = \code{0x3A}$, pesan $m = \code{0x41}$, serta $\mathrm{ipad} = \code{0x36}$, $\mathrm{opad} = \code{0x5C}$.

\begin{itemize}
    \item Kunci efektif dalam: $K \oplus \mathrm{ipad} = \code{0x3A} \oplus \code{0x36} = \code{0x0C} = 12$.
    \item Kunci efektif luar: $K \oplus \mathrm{opad} = \code{0x3A} \oplus \code{0x5C} = \code{0x66} = 102$.
    \item Lapis dalam: $h_1 = H(\code{0x0C}, \code{0x41}) = (12 + 65) \bmod 256 = 77$.
    \item Lapis luar: $\mathrm{Tag} = H(\code{0x66}, 77) = (102 + 77) \bmod 256 = 179$.
\end{itemize}
Hasilnya $\mathrm{HMAC}(K,m) = 179$.

\textbf{Perbandingan dengan $H(K\|m)$.} Untuk skema naif $H(K\|m) = (58 + 65) \bmod 256 = 123 = \code{0x7B}$. Penyerang yang tahu tag ini dan panjang $K$ dapat menghitung $H(K\|m\|\code{0x42})$ tanpa $K$, yaitu $123 + 66 = 189 \pmod{256}$; penambahan byte \code{0x42} mengubah tag secara terprediksi. Serangan \textit{length extension} inilah yang tidak mungkin dilakukan terhadap HMAC dua lapis.

\begin{obeactivity}
\textbf{Aktivitas 15.1: Eksperimen Efek Longsoran (\textit{Avalanche Effect}) pada Hash}
1. Gunakan alat bantu online (seperti CyberChef) untuk menghitung SHA-256 dari kalimat: \texttt{"Kriptografi itu Menyenangkan"}.
2. Ubah satu karakter kecil, misalnya menjadi: \texttt{"Kriptografi itu menyenangkan"} (huruf 'm' kecil).
3. Bandingkan kedua hasil hash tersebut.
4. Amati berapa banyak bit yang berubah dan diskusikan mengapa properti ini sangat penting bagi keamanan fungsi hash.
\end{obeactivity}

Langkah-langkah pada Aktivitas 15.1 dapat dijalankan langsung melalui program
pada Listing~\ref{lst:hash-avalanche}. Program itu menghitung SHA-256 dari
kedua kalimat, membandingkan digestnya, dan menghitung berapa bit yang berbeda.
Hasil yang seharusnya Anda peroleh diperlihatkan pada
Gambar~\ref{fig:longsoran-sha256}: $115$ dari $256$ bit berubah, yaitu
$44{,}92\%$ --- mendekati setengah, sebagaimana disyaratkan oleh uji longsoran
pada Subbagian~\ref{subsec:anatomi-sha256}. Bila angka yang Anda peroleh
berbeda jauh dari setengah, ada yang salah pada pesan yang di-hash, bukan pada
SHA-256.

\lstinputlisting[
    language=Python,
    caption={Efek longsoran SHA-256: perbandingan bit dua digest},
    label={lst:hash-avalanche}
]{\codefile{python/hash_avalanche.py}}

Program yang sama tersedia dalam C++ (\texttt{code/cpp/hash\_avalanche.cpp})
dan MATLAB (\texttt{code/matlab/hash\_avalanche.m}); keduanya memuat
implementasi SHA-256 sendiri karena bahasanya tidak menyediakannya secara
bawaan.

\begin{obeactivity}
\textbf{Aktivitas 15.2: Menelusuri Jejak Kompresi SHA-256}
\label{act:hash-sha256}
1. Jalankan \texttt{python code/python/hash\_uji.py} dan catat keadaan register
   $a$ dan $e$ pada putaran 1, 2, 16, 32, 48, dan 64.
2. Bandingkan digest yang dihasilkan dengan keluaran \texttt{hashlib.sha256}.
3. Ubah pesan masukan menjadi \texttt{abd} dan amati pada putaran ke berapa
   keadaan register mulai berbeda dari \texttt{abc}.
4. Periksa hasil perhitungan padding untuk tiga panjang pesan yang tercetak:
   mengapa panjang 448 bit justru memerlukan dua blok, bukan satu?
\end{obeactivity}

Berbeda dengan Aktivitas 15.1 yang memakai \texttt{hashlib} sebagai kotak hitam,
aktivitas ini membuka isi kotaknya. SHA-256 ditulis dari definisi sehingga setiap
putaran dapat diperiksa satu per satu. Bagian jadwal pesan dan fungsi kompresi
(Persamaan~\eqref{eq:sha256-jadwal} dan~\eqref{eq:sha256-kompresi}) diperlihatkan
pada Listing~\ref{lst:hash-jadwal}, sedangkan aturan padding
(Persamaan~\eqref{eq:padding-sha}) beserta pengulangan bloknya ada pada
Listing~\ref{lst:hash-padding}.

\lstinputlisting[
    language=Python,
    caption={Jadwal pesan dan fungsi kompresi SHA-256 selama 64 putaran},
    label={lst:hash-jadwal},
    firstline=60,
    lastline=90
]{\codefile{python/hash_uji.py}}

\lstinputlisting[
    language=Python,
    caption={Aturan padding Merkle--Damgård dan pengulangan blok SHA-256},
    label={lst:hash-padding},
    firstline=92,
    lastline=110
]{\codefile{python/hash_uji.py}}

Hasil menjalankan program untuk pesan \texttt{abc} memperlihatkan bahwa sesudah
satu putaran saja register $a$ sudah berubah dari \code{6a09e667} menjadi
\code{5d6aebcd}, dan sesudah 64 putaran digestnya berakhir dengan
\code{ba7816bf} --- tepat sama dengan nilai baku. Perhatikan juga bahwa peredaran
nilai pada register tidak memperlihatkan pola apa pun terhadap kesederhanaan
masukan \texttt{abc}. Itulah wujud konkret sifat \emph{diffusion} yang dibahas
pada Subbagian~\ref{subsec:konstruksi-hash}.

\begin{obeactivity}
\textbf{Aktivitas 15.3: Membuktikan Serangan Panjang-Ekstensi}
1. Jalankan program dan catat nilai $H(K \| m)$ untuk $K = \texttt{"kunci"}$ dan
   $m = \texttt{"pesan"}$.
2. Perhatikan bahwa penyerang hanya mengetahui panjang $|K \| m| = 10$ byte,
   bukan isi $K$.
3. Bandingkan tag palsu yang dihitung dari keadaan internal dengan hash yang
   dihitung \texttt{hashlib} atas pesan yang diperpanjang. Apakah keduanya sama?
4. Diskusikan: berapa operasi tambahan yang dibutuhkan penyerang dibandingkan
   dengan menghitung hash pesan biasa?
\end{obeactivity}

Serangan pada Aktivitas 15.3 memperlihatkan bahwa skema $H(K \| m)$ dapat
ditembus tanpa mengetahui kunci sama sekali. Seluruh serangan itu dijalankan oleh
fungsi pada Listing~\ref{lst:hash-ekstensi}, yang hanya melanjutkan kompresi dari
digest yang sudah diketahui. Hasil numeriknya dibahas pada
Contoh~\ref{ex:length-extension}.

\lstinputlisting[
    language=Python,
    caption={Inti serangan panjang-ekstensi: melanjutkan kompresi dari digest},
    label={lst:hash-ekstensi},
    firstline=117,
    lastline=130
]{\codefile{python/hash_uji.py}}

\begin{obeactivity}
\textbf{Aktivitas 15.4: Menguji MAC: HMAC, Vektor RFC 4231, dan Waktu-Konstan}
1. Jalankan program dan periksa apakah kelima vektor uji HMAC-SHA-256 cocok
   dengan keluaran \texttt{hmac.new}.
2. Perhatikan TC6: kuncinya 131 byte, melebihi ukuran blok 64 byte. Cabang mana
   pada Persamaan~\eqref{eq:hmac-panjang-kunci} yang dipakai?
3. Hitung sendiri batas ulang tahun untuk $n = 128$, $160$, $256$, dan $512$,
   lalu bandingkan dengan nilai yang tercetak.
4. Pada bagian terakhir, kedua fungsi pemeriksa tag menolak tag yang diubah.
   Jelaskan mengapa cara pertama tetap berbahaya meskipun jawabannya benar.
\end{obeactivity}

Listing~\ref{lst:hash-hmac} memuat HMAC yang ditulis mengikuti
Persamaan~\eqref{eq:hmac-rumus}, sedangkan Listing~\ref{lst:hash-uji} memuat
program pengujiannya: jejak kompresi, aturan padding, serangan panjang-ekstensi,
kelima vektor RFC 4231, batas ulang tahun, dan perbandingan tag secara
waktu-konstan.

\lstinputlisting[
    language=Python,
    caption={HMAC-SHA-256 dari definisi: penyesuaian panjang kunci dan dua lapis hash},
    label={lst:hash-hmac},
    firstline=136,
    lastline=146
]{\codefile{python/hash_uji.py}}

\lstinputlisting[
    language=Python,
    caption={Program pengujian: jejak kompresi, panjang-ekstensi, vektor RFC 4231,
    batas ulang tahun, dan perbandingan waktu-konstan},
    label={lst:hash-uji},
    firstline=166,
    lastline=263
]{\codefile{python/hash_uji.py}}

Sesudah aktivitas ini, mahasiswa diharapkan dapat menjawab pertanyaan yang
menjadi inti bab ini: bukan sekadar \emph{apa} fungsi hash itu, melainkan
\emph{mengapa} setiap pilihan konstruksinya diambil, dan apa yang terjadi bila
pilihan itu keliru.


\begin{obereflection}
\textbf{Latihan 15.1}
1. Mengapa \textit{Collision Resistance} lebih sulit dicapai daripada \textit{Pre-image Resistance}?
2. Apa perbedaan mendasar antara penggunaan fungsi hash untuk \textit{checksum} file dengan penggunaan hash dalam tanda tangan digital?
3. Dalam skenario apa HMAC lebih tepat digunakan dibandingkan hanya menggunakan fungsi hash biasa?
\end{obereflection}

\section{Rangkuman}
\begin{itemize}
    \item Fungsi hash memetakan pesan berukuran variabel menjadi \textit{digest} berukuran tetap dan bersifat satu arah; ia menjamin integritas, bukan kerahasiaan. Ia dibangun dari fungsi kompresi yang bekerja atas blok berukuran tetap.
    \item Properti kriptografisnya: tahan \textit{pre-image}, tahan \textit{second pre-image}, dan tahan kolisi; ketahanan kolisi hanya sebesar $2^{n/2}$ karena \textit{birthday paradox} --- pada $k = 1{,}177 \cdot 2^{n/2}$ peluang kolisi mencapai $50\%$, dan pada $k = 2^{n/2}$ mencapai $39{,}3\%$.
    \item Padding Merkle--Damgård menyisipkan bit \code{1}, lalu bit \code{0} sampai panjang bersisa 448 modulo 512, lalu panjang pesan sebagai 64 bit. Karena panjang ikut diproses sebagai data, menambahkan bit pada akhir pesan tidak dapat menghasilkan hash yang sama.
    \item Parameter konkret: MD5 128 bit/512 bit (rusak), SHA-1 160 bit/512 bit (kolisi praktis 2017), SHA-256 256 bit/512 bit/64 putaran, SHA-512 512 bit/1024 bit/80 putaran, SHA-3 berbasis sponge Keccak.
    \item MD5, SHA-1, dan SHA-2 memakai konstruksi Merkle--Damgård yang rawan \textit{length extension}; SHA-3 memakai konstruksi sponge dengan pemisahan \textit{rate} dan kapasitas, sehingga tahan terhadap serangan itu dan keamanannya tidak bergantung pada ketahanan kolisi sebuah fungsi kompresi.
    \item Konstanta SHA-2 dapat diperiksa asalnya: $K_t$ dari bagian pecahan akar pangkat tiga 64 bilangan prima pertama, dan nilai awalnya dari akar kuadrat 8 bilangan prima pertama. Cara itu mencegah penyisipan pintu belakang melalui pilihan konstanta.
    \item Serangan kolisi nyata telah menjatuhkan MD5 (Wang, 2004) dan SHA-1 (\textit{SHAttered}, 2017, sekitar $2^{63{,}1}$ operasi kompresi), serta mencapai bentuk \textit{chosen-prefix} pada keduanya. Serangan-serangan itu menempuh kelemahan struktur kompresi, bukan penelusuran menyeluruh.
    \item MAC menambahkan kunci simetris pada fungsi hash sehingga menjamin integritas sekaligus otentikasi pesan; HMAC memakai dua lapis hash dengan $\mathrm{ipad} = \code{0x36}$ dan $\mathrm{opad} = \code{0x5C}$, dan keamanannya bersandar pada sifat pseudorandom, bukan pada ketahanan kolisi.
    \item HMAC lebih kuat daripada $H(K\|m)$ karena menutup serangan \textit{length extension}, namun keduanya tidak memberikan \textit{non-repudiation}. Selain HMAC, tersedia CMAC (dua subkunci $K_1$ dan $K_2$), GMAC, dan Poly1305.
    \item AEAD seperti GCM, CCM, dan ChaCha20-Poly1305 menghasilkan cipherteks dan tag sekaligus, tetapi semuanya menuntut nonce yang unik untuk setiap pesan; pengulangan nonce pada GCM meruntuhkan kerahasiaan sekaligus memungkinkan pemalsuan.
    \item Checksum hanya mendeteksi galat acak, MAC menjamin otentikasi dengan kunci bersama, dan tanda tangan digital menjamin integritas, otentikasi, serta \textit{non-repudiation}.
    \item Pemeriksaan tag harus dilakukan secara waktu-konstan. Perbandingan yang berhenti pada byte pertama yang berbeda membocorkan panjang awalan tag yang benar, dan menurunkan biaya pemalsuan dari $2^{t}$ menjadi sekitar $t \cdot 256$.
\end{itemize}


\begin{obeassessment}
\textbf{Kuis Bab 15}
1. Jelaskan apa yang dimaksud dengan \textit{One-Way Function} dan berikan contoh penerapannya.
2. Bandingkan antara SHA-256 dan MD5 dari sisi panjang output dan tingkat keamanan.
3. Bagaimana mekanisme kerja HMAC dalam menjamin integritas sekaligus otentikasi pesan?
4. Tuliskan aturan padding Merkle--Damgård dan hitung berapa bit nol yang disisipkan untuk pesan sepanjang 1000 bit. Berapa blok yang terbentuk?
5. Mengapa ketahanan kolisi sebuah fungsi hash hanya sebesar $2^{n/2}$? Turunkan hubungan itu dari paradoks ulang tahun.
6. Sebutkan nilai awal (IV) SHA-256 dan jelaskan dari mana kedelapan nilai itu diperoleh.
7. Jelaskan mengapa skema $H(K \| m)$ tidak aman, sedangkan $H(m \| K)$ gagal karena alasan yang berbeda. Serangan apa yang menjatuhkan masing-masing?
8. Bandingkan HMAC, CMAC, dan Poly1305 dari sisi bahan dasar, kecepatan, dan lingkungan pemakaian yang cocok.
9. Apa yang dimaksud dengan AEAD? Mengapa ketiga skema AEAD yang dibahas menuntut nonce yang unik, dan apa akibat pelanggaran syarat itu pada GCM?
10. Jelaskan mengapa perbandingan tag yang berhenti pada byte pertama yang berbeda membahayakan, dan tuliskan bentuk pemeriksaan yang aman.
\end{obeassessment}
\begin{competencychecklist}
\begin{tabularx}{\linewidth}{|X|c|}
\hline
Indikator Kompetensi & Tercapai \\
\hline
Saya dapat membedakan fungsi hash dengan enkripsi & [ ] \\
\hline
Saya dapat menjelaskan tiga properti utama hash kriptografis & [ ] \\
\hline
Saya dapat menurunkan batas ulang tahun dan menerapkannya pada MD5 dan SHA-256 & [ ] \\
\hline
Saya dapat menjelaskan aturan padding Merkle--Damgård beserta akibatnya & [ ] \\
\hline
Saya dapat membandingkan algoritma MD5, SHA-1, SHA-2, dan SHA-3 & [ ] \\
\hline
Saya dapat menelusuri jejak kompresi SHA-256 dan memeriksanya dengan Python & [ ] \\
\hline
Saya dapat menjelaskan serangan panjang-ekstensi dan cara menutupnya & [ ] \\
\hline
Saya dapat menguraikan cara kerja MAC dan HMAC & [ ] \\
\hline
Saya dapat membandingkan HMAC, CMAC, GMAC, dan Poly1305 & [ ] \\
\hline
Saya dapat menjelaskan prinsip AEAD dan syarat keunikan nonce & [ ] \\
\hline
Saya dapat menganalisis risiko serangan kolisi pada fungsi hash & [ ] \\
\hline
Saya dapat menuliskan pemeriksaan tag yang tahan terhadap serangan waktu & [ ] \\
\hline
\end{tabularx}
\end{competencychecklist}


\end{document}
